\documentclass{report}

%%%%%%%%%%%%%%%%%%%%%%%%%% Загальна інформація %%%%%%%%%%%%%%%%%%%%%%%%%%
\Subject {Технології графічного процесінгу}
\LabTitle{Прискорені обчислення для Python з Numba}
\LabReport{Практична робота \#4}


%----------------------------------------------------
%	  Персональна інформація виконавця завдання
%----------------------------------------------------
\Done{Виконав:}  % Поставте правильне закінчення


\Surname{Ковальов}  % Вкажіть своє прізвище
\Name{Олександр}			% та ім'я
\Group {ІМ-51мн}     % група в якій навчаєтесь
\YearOfStudying {1} % курс навчання


%%%%%%%%%%%%%%%%%%%%%%%%%%	 ПОЧАТОК ЗВІТУ	 %%%%%%%%%%%%%%%%%%%%%%%%
\startDocument

\section{Мета}
Навчитися використовувати Numba для компiляцiї та прискорення Python-функцiй на CPU i GPU, писати власнi векторизованi функцiї з виконанням на GPU, а також керувати передачею даних мiж CPU i GPU, щоб уникнути зайвих накладних витрат.

\section{Завдання}

Вiдкрийте стартовий код завдання за посиланням: \url{colab.research.google.com/github/YKochura/ac-kpi/blob/main/labs/pr4\_python\_numba.ipynb}

Реалiзуйте три версiї функцiї \texttt{create\_hidden\_layer}: на чистому NumPy (стартовий код завдання), з Numba CPU (\texttt{@vectorize}) та з Numba GPU (\texttt{@vectorize(target=’cuda’)}). Для кожної реалiзацiї вимiряйте середнiй час виконання (рекомендується попередньо виконати JIT-прогрiвання). Numba при першому виклику функцiї компiлює Python-код у машинний код пiд конкретний тип даних -- це займає десятки мiлiсекунд i не вiдображає реальну швидкiсть обчислень. Усi наступнi виклики вже використовують скомпiльований код i працюють швидше. 

В GPU-версiї данi мiж CPU i GPU передавайте лише один раз, а не пiсля кожної операцiї. 

Результати оформiть у виглядi протоколу: таблицю замiрiв для рiзної кiлькостi нейронiв \texttt{n = 1M} i \texttt{n = 100M}, порiвняння продуктивностi та висновки щодо отриманого прискорення i ролi оптимiзацiї передачi даних.

\section*{Протокол виконання лабораторної роботи}

%% Section: Benchmark Table
\subsection*{Таблиця замірів продуктивності}

\begin{table}[h]
	\centering
	\caption{Результати вимірювання часу виконання функції \texttt{create\_hidden\_layer}}
	\vspace{0.2cm}
	\begin{tabular}{|l|c|c|}
		\hline
		\textbf{Версія реалізації} & \textbf{Час виконання ($n = 10^6$)} & \textbf{Час виконання ($n = 10^8$)} \\
		\hline
		Чистий NumPy & 9.74 ms & 1.87 s \\
		\hline
		Numba CPU & 8.89 ms & 970 ms \\
		\hline
		Numba GPU & 3.88 ms & 295 ms \\
		\hline
	\end{tabular}
\end{table}

%% Section: Performance Comparison
\subsection*{Порівняння продуктивності}

Аналіз отриманих результатів вимірювання часу виконання яскраво демонструє переваги використання JIT-компіляції та апаратного прискорення. Базова реалізація на чистому NumPy показала найнижчу продуктивність під час роботи з великим масивом даних розміром у сто мільйонів елементів, витративши 1.87 секунди. Це зумовлено тим, що під час виконання векторизованих математичних операцій створюються нові тимчасові масиви в оперативній пам'яті для кожного проміжного кроку. Використання JIT-компілятора Numba для центрального процесора дозволило скоротити час виконання до 970 мілісекунд, забезпечивши прискорення майже вдвічі. Декоратор \texttt{vectorize} зливає всі дрібні операції в один оптимізований цикл на рівні машинного коду, усуваючи необхідність створення проміжних масивів. Реалізація з використанням Numba CUDA для графічного процесора продемонструвала найвищу ефективність, виконавши завдання за 295 мілісекунд, що у понад шість разів швидше за базовий підхід. На меншому обсязі даних в один мільйон елементів різниця між процесорними підходами мінімальна через накладні витрати на виклик функцій, проте GPU все одно зберігає лідерство із результатом 3.88 мілісекунди.

%% Section: Conclusions
\subsection*{Висновки щодо отриманого прискорення та ролі оптимізації передачі даних}

Застосування графічного процесора для векторизованих обчислень забезпечує суттєвий приріст продуктивності завдяки масивній паралелізації, що особливо помітно на великих обсягах даних. Вирішальну роль у досягненні такого високого показника відіграла оптимізація передачі даних між оперативною пам'яттю хоста та відеопам'яттю пристрою. Пропускна здатність шини PCIe є фізичним обмеженням, яке може повністю нівелювати обчислювальну потужність відеокарти при нераціональному використанні. Одноразове завантаження початкових масивів у пам'ять відеокарти перед початком обчислень, виконання всіх математичних перетворень локально на пристрої та одноразове повернення фінального результату дозволили уникнути критичних накладних витрат. Отже, ефективне управління пам'яттю є настільки ж важливим етапом у технологіях графічного процесінгу, як і безпосереднє розпаралелювання обчислень.

\newpage

\section{Лістинг.}

\begin{lstlisting}[language=Python, caption={Код.}] 
	import numpy as np
	import math
	from numba import vectorize, cuda
	
	\# Set initial number of neurons for the first test
	n = 1\_000\_000
	
	\# Generate initial data
	greyscales = np.floor(np.random.uniform(0, 255, n).astype(np.float32))
	weights = np.random.normal(.5, .1, n).astype(np.float32)
	
	\# Pack arguments for convenience
	arguments = \{
	``greyscales'': greyscales,
	``weights'': weights
	\}
	
	\# Pure NumPy implementation
	
	def normalize\_np(grayscales):
	return grayscales / 255.0
	
	def weigh\_np(values, weights):
	return values * weights
	
	def activate\_np(values):
	return (np.exp(values) - np.exp(-values)) / (np.exp(values) + np.exp(-values))
	
	def create\_hidden\_layer\_numpy(greyscales, weights):
	normalized = normalize\_np(greyscales)
	weighted = weigh\_np(normalized, weights)
	return activate\_np(weighted)
	
	\# Numba CPU implementation
	
	@vectorize(['float32(float32)'], target='cpu')
	def normalize\_cpu(grayscales):
	return grayscales / 255.0
	
	@vectorize(['float32(float32, float32)'], target='cpu')
	def weigh\_cpu(values, weights):
	return values * weights
	
	@vectorize(['float32(float32)'], target='cpu')
	def activate\_cpu(values):
	\# Use math.exp instead of numpy.exp for Numba compatibility
	return (math.exp(values) - math.exp(-values)) / (math.exp(values) + math.exp(-values))
	
	def create\_hidden\_layer\_cpu(greyscales, weights):
	normalized = normalize\_cpu(greyscales)
	weighted = weigh\_cpu(normalized, weights)
	return activate\_cpu(weighted)
	
	\# Numba GPU implementation with optimized data transfer
	
	@vectorize(['float32(float32)'], target='cuda')
	def normalize\_gpu(grayscales):
	return grayscales / 255.0
	
	@vectorize(['float32(float32, float32)'], target='cuda')
	def weigh\_gpu(values, weights):
	return values * weights
	
	@vectorize(['float32(float32)'], target='cuda')
	def activate\_gpu(values):
	return (math.exp(values) - math.exp(-values)) / (math.exp(values) + math.exp(-values))
	
	def create\_hidden\_layer\_gpu(greyscales, weights):
	\# Transfer arrays to device memory once to avoid PCIe bottleneck
	d\_greyscales = cuda.to\_device(greyscales)
	d\_weights = cuda.to\_device(weights)
	
	\# Perform all calculations on the device
	d\_normalized = normalize\_gpu(d\_greyscales)
	d\_weighted = weigh\_gpu(d\_normalized, d\_weights)
	d\_activated = activate\_gpu(d\_weighted)
	
	\# Copy the final result back to host memory
	return d\_activated.copy\_to\_host()
	
	\# JIT Warm-up sequence
	\# Execute functions with a small subset to trigger compilation
	
	warmup\_g = greyscales[:10]
	warmup\_w = weights[:10]
	
	\_ = create\_hidden\_layer\_numpy(warmup\_g, warmup\_w)
	\_ = create\_hidden\_layer\_cpu(warmup\_g, warmup\_w)
	\_ = create\_hidden\_layer\_gpu(warmup\_g, warmup\_w)
	
	print(``Warm-up complete. Functions are compiled.'')
	
	\# Benchmark for n = 1M
	
	print(``NumPy version (1M):'')
	\%timeit create\_hidden\_layer\_numpy(**arguments)
	
	print(``\textbackslash{}nNumba CPU version (1M):'')
	\%timeit create\_hidden\_layer\_cpu(**arguments)
	
	print(``\textbackslash{}nNumba GPU version (1M):'')
	\%timeit create\_hidden\_layer\_gpu(**arguments)
	
	\# Generate data and benchmark for n = 100M
	
	n\_large = 100\_000\_000
	
	greyscales\_large = np.floor(np.random.uniform(0, 255, n\_large).astype(np.float32))
	weights\_large = np.random.normal(.5, .1, n\_large).astype(np.float32)
	
	arguments\_large = \{
	``greyscales'': greyscales\_large,
	``weights'': weights\_large
	\}
	
	print(``NumPy version (100M):'')
	\%timeit create\_hidden\_layer\_numpy(**arguments\_large)
	
	print(``\textbackslash{}nNumba CPU version (100M):'')
	\%timeit create\_hidden\_layer\_cpu(**arguments\_large)
	
	print(``\textbackslash{}nNumba GPU version (100M):'')
	\%timeit create\_hidden\_layer\_gpu(**arguments\_large)
\end{lstlisting}

\end{document}
